\documentclass[11pt]{article}
\usepackage[margin=0.75in]{geometry}
\usepackage{times}
\usepackage{array}
\usepackage{hyperref}

\title{Rebuttal to Reviews}
\author{Anonymous Authors}
\date{}

\begin{document}
\maketitle

\section*{Summary of Changes}
\begin{itemize}
  \item \textbf{Reviewer 2, W1 (soundness):} Added ablation study on X; results show $\Delta=+Y\%$ (Table 2).
  \item \textbf{Reviewer 3, W2 (novelty):} Clarified positioning vs. recent work Z in \S3.1.
  \item \textbf{Reviewer 1, W3 (clarity):} Rewrote \S4.2 for self-contained explanation.
\end{itemize}

\section*{Response to Reviewer 1}
\textbf{W3 (clarity).} Thank you. We rewrote \S4.2 to make the derivation self-contained; the revised text now states the assumption explicitly and adds a worked example.

\section*{Response to Reviewer 2}
\textbf{W1 (soundness).} We agree and added the missing ablation on X. The new experiment (Table 2) shows our method gains $+Y\%$ over the strongest baseline, directly addressing your concern. \\
\textbf{W4 (baseline).} We now compare against the concurrent SOTA B (cited in Related Work) and report the gap in Table 3.

\section*{Response to Reviewer 3}
\textbf{W2 (novelty).} As discussed in \S3.1, our contribution is orthogonal to Z: Z targets setting A, while we target setting B. We added a paragraph distinguishing the two and cite Z appropriately.

\section*{Response to the Area Chair}
All three primary concerns have been directly addressed: R2's soundness critique is answered with new experiments, R3's novelty concern with a positioning clarification, and R1's clarity point with a rewrite. We believe the revision substantially strengthens the paper.

\end{document}
